\phantomsection
\section*{修订历史}\label{revision-history}
\lhead{\it{修订历史}}
\rhead{\Navigation}
\addcontentsline{toc}{chapter}{修订历史}

\begin{longtable}[l]{ | m{2cm} | m{1.5cm} | m{12cm} | }\hline

    \rowcolor{lightgray}
    日期 & 版本 & \hypertarget{Release notes}{发布说明}\\\hline
    \endfirsthead

    \multicolumn{3}{c}{\bfseries 接上页}\\\hline
    \rowcolor{lightgray}
    日期 & 版本 & 发布说明\\\hline
    \endhead

    \multicolumn{3}{r}{\textbf{见下页}}\\
    \endfoot

    \endlastfoot

    % table contents
    2026-03-20 & v1.2 &
    更新以下章节：
    \begin{itemize}
        \item 章节 \ref{mod:ulp} \textit{\nameref{mod:ulp}}：更新 LP CPU 和 HP CPU 原子访问的限制 %DOC-8741
        \item 章节 \ref{mod:lowpowm} \textit{\nameref{mod:lowpowm}}：新增寄存器 \hyperref[regdesc:PMUPOWERVDDSPICNTLREG]{PMU\_POWER\_VDD\_SPI\_CNTL\_REG} 的描述，更正错别字 %DOC-12323 DOC-14058 DOC-14057
        \item 章节 \ref{mod:digsig} \textit{\nameref{mod:digsig}}：将模块从数字签名 (DS) 更名为 RSA 数字签名外设 (RSA\_DS) %DOC-13833
        \item 章节 \ref{mod:i2c} \textit{\nameref{mod:i2c}}：更新从机启动方式的描述 %DOC-12530
        \item 章节 \ref{mod:rmt} \textit{\nameref{mod:rmt}}：更新 RMT 时钟频率限制相关公式 %DOC-13581
    \end{itemize}
    \\\hline

    2025-06-18 & v1.1 &
    更新以下章节：
    \begin{itemize}
        \item 章节 \ref{mod:dma} \textit{\nameref{mod:dma}}：新增 GDMA\_OUTFIFO\_OVF\_CH\regindex{n}\_INT、GDMA\_OUTFIFO\_UDF\_CH\regindex{n}\_INT、GDMA\_INFIFO\_OVF\_CH\regindex{n}\_INT 和 GDMA\_INFIFO\_UDF\_CH\regindex{n}\_INT 中断的描述 %DOC-5353
        \item 章节 \ref{mod:resclk} \textit{\nameref{mod:resclk}}：根据 \href{}{AR2024-011 关于 ESP32-C6 芯片删除 RC32K 时钟的说明}，移除 RC32K\_CLK；更新关于 RC\_SLOW\_CLK 时钟的描述 %DOC-6671 %DOC-9387
        \item 章节 \ref{mod:intmtrx} \textit{\nameref{mod:intmtrx}}：移除 CACHE\_INTR 中断源 %DOC-8488
        \item 章节 \ref{mod:lowpowm} \textit{\nameref{mod:lowpowm}}：修正表 \ref{tab:lowpowm-wakeup-source}，LP CPU 也可作为 Deep-sleep 的唤醒源；根据 \href{}{AR2024-011 关于 ESP32-C6 芯片删除 RC32K 时钟的说明}，移除 RC32K\_CLK；增加 RTC\_EVT\_TICK 的事件描述；增加寄存器 \hyperref[fielddesc:PMULPCPUSLPSTALLWAIT]{PMU\_LP\_CPU\_SLP\_STALL\_WAIT}、\hyperref[fielddesc:LPAONFORCEDOWNLOADBOOT]{LP\_AON\_FORCE\_DOWNLOAD\_BOOT}、\hyperref[regdesc:LPAONGPIOMUXREG]{LP\_AON\_GPIO\_MUX\_REG}、\hyperref[regdesc:LPAONGPIOHOLD0REG]{LP\_AON\_GPIO\_HOLD0\_REG}、\hyperref[regdesc:LPAONLPBUSREG]{LP\_AON\_LPBUS\_REG}、\hyperref[regdesc:PMUIMMPADHOLDALLREG]{PMU\_IMM\_PAD\_HOLD\_ALL\_REG} 的描述 %DOC-6671 %DOC-10466 %DOC-10808 %DOC-10554 %DOC-10553
        \item 章节 \ref{mod:ecc} \nameref{mod:ecc}：将寄存器字段前缀从 ECC 更新为 ECC\_MULT %DOC-10131
        \item 章节 \ref{mod:uart} \textit{\nameref{mod:uart}}：更新 FIFO 相关描述，澄清发送块和接收块的 RAM 存储空间以及 FIFO 使用方式；新增写 \hyperref[fielddesc:UARTRXFIFORDBYTE]{UART\_RXFIFO\_RD\_BYTE} 的说明 %DOC-8553, DOC-9040
        \item 章节 \ref{mod:sensor} \textit{\nameref{mod:sensor}}：将 ADC 滤波公式中的 ``-0.5" 更正为 ``+0.5"；新增温度传感器的 ETM 任务描述 %DOC-10226 %DOC-10296
    \end{itemize}
    \\\hline
    2024-06-24        & v1.0 &
    新增以下章节：
    \begin{itemize}
        \item 章节 \ref{mod:lowpowm} \nameref{mod:lowpowm}
    \end{itemize}
    更新以下章节：
    \begin{itemize}
        \item 章节 \ref{mod:ulp} \textit{\nameref{mod:ulp}}：更新 mie 和 mip 寄存器描述 %DOC-8097
        \item 章节 \ref{mod:resclk} \textit{\nameref{mod:resclk}}：新增 PMU 控制的时钟门控信息 % DOC-8149
        \item 章节 \ref{mod:spi} \textit{\nameref{mod:spi}}：更新 GP-SPI2 用作从机时支持的输入时钟频率 %DOC-8034
        \item 将 ``LP Timer" 全局更新为 ``RTC Timer" %DOC-7121
    \end{itemize}
    \\\hline
    2024-05-10        & v0.4 &
    更新以下章节：
    \begin{itemize}
        \item 章节 \ref{mod:ulp} \textit{\nameref{mod:ulp}}：更新若干 PMU 寄存器名称 %DOC-7713
        \item 章节 \ref{mod:dma} \textit{\nameref{mod:dma}}：更新 suc\_eof 和 EOF 标志的相关描述 %DOC-5478
        \item 章节 \ref{mod:sysmem} \textit{\nameref{mod:sysmem}}：在表 \ref{tab:sysmem-base-address} 中新增 片外存储器加密与解密 的基地址；更新小节 \ref{para:sysmem-internal-memory} 中的措辞 %DOC-6666 % DOC-7280
        \item 章节 \ref{mod:efuse} \textit{\nameref{mod:efuse}}：更新EFUSE\_CRYPT\_DPA\_ENABLE 域的位置 %DOC-6338
        \item 章节 \ref{mod:iomuxgpio} \textit{\nameref{mod:iomuxgpio}}：更新 \ref{PadHold} 中的描述，删除 GPIO\_PCPU\_NMI\_INT\_REG 寄存器及相关信息 %DOC-5892 %DOC-6031 %DOC-7326
        \item 章节 \ref{mod:resclk} \textit{\nameref{mod:resclk}}：在 \ref{sec:periperal-clock-configuration} 小节下方新增一条说明；更新对 \hyperlink{fielddesc:LPCLKRSTSLOWCLKSEL}{LP\_CLKRST\_SLOW\_CLK\_SEL} 的描述；新增 \hyperref[fielddesc:PCRFOSCTICKNUM]{PCR\_FOSC\_TICK\_NUM} 字段的描述 % DOC-5852 DOC-6867 DOC-6670
        \item 章节 \ref{mod:bootctrl} \textit{\nameref{mod:bootctrl}}：新增 EFUSE\_DIS\_FORCE\_DOWNLOAD 和 EFUSE\_DIS\_DOWNLOAD\_MODE 对芯片 Boot 控制的细节描述 %DOC-5912
        \item 章节 \ref{mod:intmtrx} \textit{\nameref{mod:intmtrx}}：移除 TG0\_T1\_INTR 和 TG1\_T1\_INTR 中断源，这两个中断源不存在；删除 INTMTX\_CORE0\_GPIO\_INTERRUPT\_PRO\_NMI\_MAP\_REG 寄存器及相关信息；移除 PAU\_INTR 中断源 %DOC-6032；TRMC6-42194
        \item 章节 \ref{mod:evntaskmatrix} \textit{\nameref{mod:evntaskmatrix}}：将支持事件任务矩阵的外设由 RTC 看门狗定时器改为 RTC 定时器 %DOC-5151
        \item 章节 \ref{mod:uart} \textit{\nameref{mod:uart}}：更新 RAM 大小相关描述和产生 wake\_up 所需的上升沿个数，更新表 \ref{tab:uart-feature-distinction} 中 LP\_UART 的时钟源 %DOC-7076, DOC-6542
        \item 章节 \ref{mod:ledpwm} \textit{\nameref{mod:ledpwm}}：更新表 \ref{tab:ledc-common-freq-res} 中的 RC\_FAST\_CLK 的频率和最低精度  %DOC-7304
        \item 章节 \ref{mod:parlio} \textit{\nameref{mod:parlio}}：修改小节 \ref{parlio-lcd-data-trans} \textit{\nameref{parlio-lcd-data-trans}} 中的示例
    \end{itemize}
    新增章节 \hyperref[interrupt-config-registers]{\textit{中断配置寄存器}}
    \\\hline
    2023-06-07        & v0.3 &
    新增以下章节：
    \begin{itemize}
        \item 章节 \ref{mod:ulp} \textit{\nameref{mod:ulp}}
    \end{itemize}
    更新以下章节：
    \begin{itemize}
        \item 章节 \ref{mod:riscvcpu} \textit{\nameref{mod:riscvcpu}}：更正中断 ID %DOC-5390
        \item 章节 \ref{mod:efuse} \textit{\nameref{mod:efuse}}：更新EFUSE\_SEC\_DPA\_LEVEL 寄存器的描述 %DOC-5323
        \item 章节 \ref{mod:resclk} \textit{\nameref{mod:resclk}}：更新寄存器部分对复位字段的描述 %DOC-5214
        \item 章节 \ref{mod:intmtrx} \textit{\nameref{mod:intmtrx}}：添加 LP\_PERI\_TIMEOUT\_INTR 中断源 %DOC-5308
        \item 章节 \ref{mod:permctrl} \textit{\nameref{mod:permctrl}}：添加 \ref{fig:pmp_apm}，并更新相关描述 %DOC-5144
        \item 章节 \ref{mod:spi} \textit{\nameref{mod:spi}}：更新时钟相关信息 %DOC-4719
        \item 章节 \ref{mod:sensor} \textit{\nameref{mod:sensor}}：更新 \ref{sec:sensor-etm-func} \textit{\nameref{sec:sensor-etm-func}} 小节
    \end{itemize}
    新增章节 \hyperref[programming-reserved-field] {\textit{如何配置寄存器的保留域}}
    \\\hline
    2023-04-20        & v0.2 &
    新增以下章节：
    \begin{itemize}
        \item 章节 \ref{mod:riscvcpu} \textit{\nameref{mod:riscvcpu}}
        \item 章节 \ref{mod:riscvtracenc} \textit{\nameref{mod:riscvtracenc}}
        \item 章节 \ref{mod:sdioslave} \textit{\nameref{mod:sdioslave}}
        \item 章节 \ref{mod:parlio} \textit{\nameref{mod:parlio}}
    \end{itemize}
    更新以下章节：
    \begin{itemize}
        \item 章节 \ref{mod:dma} \textit{\nameref{mod:dma}}：更新 GDMA\_IN\_SUC\_EOF\_CH\regindex{n}\_INT 中断和 GDMA\_INLINK\_DSCR\_ADDR\_CH\regindex{n} 字段的描述 % DOC-4730
        \item 章节 \ref{mod:efuse} \textit{\nameref{mod:efuse}}：增加烧写 XTS-AES 密钥时需要注意的事项
        \item 章节 \ref{mod:resclk} \textit{\nameref{mod:resclk}}: RC\_FAST\_CLK 更新为 17.5 MHz, RC\_SLOW\_CLK 更新为 136 kHZ. %IDFH-9325
        \item 章节 \ref{mod:intmtrx} \textit{\nameref{mod:intmtrx}}：更新几个寄存器的名称
        \item 章节 \ref{mod:evntaskmatrix} \textit{\nameref{mod:evntaskmatrix}}：更新 ETM 通道的完整配置过程
        \item 章节 \ref{mod:dbgassist} \textit{\nameref{mod:dbgassist}}：增加ASSIST\_DEBUG\_CLOCK\_GATE\_REG 和 MEM\_MONITOR\_CLOCK\_GATE\_REG 两个寄存器的描述；更新表 \ref{tab:dma_source}
        \item 章节 \ref{mod:rsa} \textit{\nameref{mod:rsa}}、\ref{mod:sha} \textit{\nameref{mod:sha}} 和 \ref{mod:aes} \textit{\nameref{mod:aes}}：更新使能各加速器的寄存器%doc-4669
        \item 章节 \ref{mod:uart} \textit{\nameref{mod:uart}}：更新停止位最大位数和相关描述 %DOC-4826
        \item 章节 \ref{mod:sensor} \textit{\nameref{mod:sensor}}：更新图 \ref{fig:sensor-apb-saradc-sar-patt-tab1-reg-pattern-tab-entry0-entry3} 和图 \ref{fig:sensor-apb-saradc-sar-patt-tab2-reg-pattern-tab-entry4-entry7} 的样式表编号
    \end{itemize}
    更新 \hyperref[glossary]{\textit{术语}} 小节中的外设相关词汇并新增寄存器访问类型“varies”
    \\\hline
    2023-01-13 & v0.1 & 首次发布\\\hline
\end{longtable}
